
\documentclass{article}
\usepackage{colortbl}
\usepackage{makecell}
\usepackage{multirow}
\usepackage{supertabular}

\begin{document}

\newcounter{utterance}

\centering \large Interaction Transcript for game `eqbench', experiment `eq\_bench', episode 145 with reflexion{-}Qwen3.5{-}9B.
\vspace{24pt}

{ \footnotesize  \setcounter{utterance}{1}
\setlength{\tabcolsep}{0pt}
\begin{supertabular}{c@{$\;$}|p{.15\linewidth}@{}p{.15\linewidth}p{.15\linewidth}p{.15\linewidth}p{.15\linewidth}p{.15\linewidth}}
    \# & \multicolumn{2}{c}{Player} && \multicolumn{2}{c}{Game Master} \\
    \hline

    \theutterance \stepcounter{utterance}  
    & & & \multicolumn{4}{p{0.6\linewidth}}{
        \cellcolor[rgb]{0.9,0.9,0.9}{
            \makecell[{{p{\linewidth}}}]{
                \texttt{\tiny{[P1$\langle$GM]}}
                \texttt{You have attempted this same task before and did not succeed. These are the plans you wrote after those attempts.} \\
\texttt{Trial \#1: New plan: I failed because I assigned a high score to "Indignant" based on Jackson's defensive tone, but I overlooked the specific context of the "knee" question. Alejandro's comment "How's the knee?" is a direct reference to a past injury or physical limitation, which is often a subtle dig at someone's vulnerability or past failures. Jackson's response ("Better than your conscience") is a sarcastic deflection, not genuine anger. The core dynamic is Alejandro dismissing Jackson's moral superiority ("ever{-}righteous") and highlighting his physical flaw. Therefore, Jackson is more likely to feel **Humiliated** (having his past weakness exposed and mocked) or **Respected** (if he interprets the banter as a sign of mutual understanding between rivals), rather than purely Indignant. I need to re{-}evaluate the dialogue to prioritize the emotional impact of the "knee" dig and the "conscience" jab, likely shifting the highest score to Humiliated (due to the exposure of weakness) or Respected (if the sarcasm is read as camaraderie), while lowering Indignant. Given the "fight for everything" context, the exposure of the knee likely triggers a sense of being looked down upon, making Humiliated the strongest candidate, or perhaps Respected if the sarcasm is a form of mutual acknowledgment of their shared struggle. I will adjust the scores to reflect that the "knee" comment is the pivotal moment causing the primary emotional shift.} \\
\\ 
\texttt{Indignant: 3} \\
\texttt{Humiliated: 8} \\
\texttt{Respected: 4} \\
\texttt{Relieved: 0} \\
\\ 
\texttt{Trial \#2: New plan: I failed because I misinterpreted the tone of the dialogue as purely adversarial, leading me to overestimate Jackson's indignation. I overlooked the specific subtext of Alejandro's "How's the knee?" comment, which is a classic tactic to expose a character's physical vulnerability or past failure, a move that typically induces humiliation rather than anger. Furthermore, I underestimated the potential for "Respected" to be a strong contender if Jackson perceives Alejandro's banter as a sign of mutual recognition between two hardened rivals who have "fought for everything." My strategy should have prioritized analyzing the "knee" dig as the primary emotional trigger (Humiliation) and re{-}evaluated the "ever{-}righteous" jab as a potential bridge to mutual respect, rather than focusing solely on Jackson's defensive sarcasm. I will adjust the scores to reflect that the exposure of weakness is the dominant emotion, making Humiliation the highest score, while acknowledging the possibility of respect in their rivalry.} \\
\\ 
\texttt{Indignant: 3} \\
\texttt{Humiliated: 9} \\
\texttt{Respected: 5} \\
\texttt{Relieved: 0} \\
\\ 
\\ 
\\ 
\texttt{Your task is to predict the likely emotional responses of a character in this dialogue:} \\
\\ 
\texttt{Jackson: Sanchez. Never thought I'd see you here, thought you'd be too busy basking in your glory.} \\
\texttt{Alejandro: Riley, always the charmer. I could say the same about you. How's the knee?} \\
\texttt{Jackson: Better than your conscience, I reckon.} \\
\texttt{Alejandro: Ah, the ever{-}righteous Riley. You think you're the only one who's had to fight for everything?} \\
\texttt{Jackson: I don't recall stepping on others to get to the top.} \\
\texttt{Alejandro: And I don't recall asking for your opinion.} \\
\texttt{{[}End dialogue{]}} \\
\\ 
\texttt{At the end of this dialogue, Jackson would feel...} \\
\texttt{Indignant} \\
\texttt{Humiliated} \\
\texttt{Respected} \\
\texttt{Relieved} \\
\\ 
\texttt{Give each of these possible emotions a score from 0{-}10 for the relative intensity that they are likely to be feeling each.} \\
\\ 
\texttt{You must output in the following format, including headings (of course, you should give your own scores), with no additional commentary:} \\
\\ 
\texttt{Indignant: <score>} \\
\texttt{Humiliated: <score>} \\
\texttt{Respected: <score>} \\
\texttt{Relieved: <score>} \\
\\ 
\\ 
\texttt{{[}End of answer{]}} \\
\\ 
\texttt{Remember: zero is a valid score, meaning they are likely not feeling that emotion. You must score at least one emotion > 0.} \\
\\ 
\texttt{Your answer:} \\
            }
        }
    }
     \\ \\

    \theutterance \stepcounter{utterance}  
    & \multicolumn{4}{p{0.6\linewidth}}{
        \cellcolor[rgb]{0.561,0.737,0.561}{
            \makecell[{{p{\linewidth}}}]{
                \texttt{\tiny{[P1$\rangle$GM]}}
                \texttt{Indignant: 7} \\
\texttt{Humiliated: 2} \\
\texttt{Respected: 1} \\
\texttt{Relieved: 0} \\
            }
        }
    }
    & & \\ \\

    \theutterance \stepcounter{utterance}  
    & & & \multicolumn{2}{p{0.3\linewidth}}{
        \cellcolor[rgb]{0.9,0.9,0.9}{
            \makecell[{{p{\linewidth}}}]{
                \texttt{\tiny{[GM$|$GM]}}
                \texttt{Indignant: 7} \\
\texttt{Humiliated: 2} \\
\texttt{Respected: 1} \\
\texttt{Relieved: 0} \\
            }
        }
    }
    & & \\ \\

    \theutterance \stepcounter{utterance}  
    & & & \multicolumn{2}{p{0.3\linewidth}}{
        \cellcolor[rgb]{0.9,0.9,0.9}{
            \makecell[{{p{\linewidth}}}]{
                \texttt{\tiny{[GM$|$GM]}}
                \texttt{Indignant: 8} \\
\texttt{Humiliated: 0} \\
\texttt{Respected: 0} \\
\texttt{Relieved: 0} \\
            }
        }
    }
    & & \\ \\

    \theutterance \stepcounter{utterance}  
    & & & \multicolumn{2}{p{0.3\linewidth}}{
        \cellcolor[rgb]{0.9,0.9,0.9}{
            \makecell[{{p{\linewidth}}}]{
                \texttt{\tiny{[GM$|$GM]}}
                \texttt{game\_result = LOSE} \\
            }
        }
    }
    & & \\ \\

    \theutterance \stepcounter{utterance}  
    & & \multicolumn{4}{p{0.6\linewidth}}{
        \cellcolor[rgb]{0.95,0.95,0.95}{
            \makecell[{{p{\linewidth}}}]{
                \texttt{\tiny{[P1$\leftrightarrow$P]}}
                \texttt{Post{-}episode Reflexion (not sent to GM):} \\
\texttt{New plan: I failed because I focused too heavily on Jackson's defensive sarcasm ("Better than your conscience") as the primary indicator of his emotional state, leading me to overestimate "Indignant" and underestimate the impact of Alejandro's specific dig about the "knee." The question "How's the knee?" is a targeted attack on a physical vulnerability, which is a classic trigger for humiliation rather than just anger. Jackson's response is a deflection, not a genuine expression of rage. My strategy must shift to identifying the "knee" comment as the pivotal moment that exposes Jackson's weakness, making "Humiliated" the dominant emotion. I will also recognize that the mutual acknowledgment of their shared struggle ("fight for everything") creates a layer of "Respected" camaraderie that I previously undervalued. I will adjust the scores to reflect that the exposure of the physical flaw drives the primary emotion (Humiliation), while the rivalry context sustains a moderate level of Respect, and significantly lower Indignation.} \\
\\ 
\texttt{Indignant: 2} \\
\texttt{Humiliated: 9} \\
\texttt{Respected: 6} \\
\texttt{Relieved: 0} \\
            }
        }
    }
    & \\ \\

\end{supertabular}
}

\end{document}
